% This file is part of the modular source for Stratified Relational Workspace Theory.
% !TeX root = ../Stratified_Relational_Workspace_Theory.tex
\theorypart{XXI}{Empirical Certification and Reporting Protocol}

\topic{74}{Required Experiments}

A workspace claim is not certified by a single benchmark score. The complete
protocol is organized as a sequence of declarations, held-out measurements,
implemented interventions, matched nulls, specification sweeps, identification
checks, and gate decisions. Every reported
result must identify the claim level it supports. Measurements that were not
performed remain unmeasured rather than being inferred from neighboring gates.

\subtopic{74.1 Claim declaration before evaluation}

Before the certification split is opened, record:
\begin{enumerate}
\item the full reference specification
 \(\varsigma_0=(\mathcal D,Y^+,\ell,\tau,\mathcal I,\mathcal N,\Theta)\);
\item the admissible decomposition family \(\mathfrak D\) and specification
 family \(\mathfrak S\);
\item the target trajectory \(Y^+\), loss, horizon, and environment family;
\item the candidate structural workspace extractor;
\item the goal family \(\mathcal G\), utility normalization, and feasible
 action set;
\item the intervention families \(\mathcal I_W,\mathcal I_G,\mathcal I_A,
 \mathcal I_D,\mathcal I_R\);
\item the null panel and the local distributions each null must preserve;
\item the required certificate level and every numerical threshold;
\item the formation, calibration, certification, and stress-test splits;
\item the uncertainty procedure and resampling unit; and
\item the stopping and model-selection rules;
\item the experiment family and model class used for identifiability; and
\item the complete specification-selection history and multiplicity rule.
\end{enumerate}
Changing any of these after observing certification outcomes creates a new
analysis that must be reported as exploratory or evaluated on a fresh split.

\subtopic{74.2 Data and intervention partition}

Use at least four logically distinct partitions:
\[
\boxed{
\mathcal D
=
\mathcal D_{\mathrm{form}}
\sqcup
\mathcal D_{\mathrm{cal}}
\sqcup
\mathcal D_{\mathrm{cert}}
\sqcup
\mathcal D_{\mathrm{stress}}.
}
\]
Their roles are:
\begin{description}
\item[Formation.] Learn membership, transfer maps, workspace channels,
 boundaries, policies, and repair models.
\item[Calibration.] Select thresholds, intervention amplitudes, and
 uncertainty procedures without making the final pass decision.
\item[Certification.] Estimate the declared gates once under the frozen
 analysis plan.
\item[Stress test.] Evaluate distribution shift, repeated damage, replacement,
 and long-horizon persistence.
\end{description}
When independent partitions are impossible, the report must state the precise
cross-fitting or nested-resampling procedure and preserve formation leakage
controls.

\subtopic{74.3 Structural tests}

Report the complete low-energy spectrum relevant to the selected threshold:
\[
\spec(L_H),
\qquad
 d_\varepsilon,
\qquad
 g_\varepsilon,
\qquad
 \delta_{\mathrm{coh}}.
\]
At minimum include:
\begin{itemize}
\item the construction and normalization of \(D_H\);
\item edge weights and estimated transfer uncertainties;
\item the selected eigenspace dimension and threshold rule;
\item the spectral gap with uncertainty or perturbation bound;
\item stability under bootstrap, seed, and representation replacement; and
\item the exact or approximate kernel residual.
\end{itemize}

\subtopic{74.4 Transfer and operation tests}

Report
\[
\eta_{\mathrm{int}},
\qquad
\delta_e^{\mathrm{tr}},
\qquad
\Omega_\gamma,
\qquad
\Lambda_\gamma,
\]
together with:
\begin{itemize}
\item the paths and cycle basis used;
\item path-length dependence of leakage;
\item the gauge-alignment procedure;
\item held-out transfer-prediction error;
\item operation-intertwining defects; and
\item the approximate closure bound from Theorem~\(T_{10}\).
\end{itemize}

\subtopic{74.5 Predictive alignment and coalition tests}

Report
\[
\delta_{\mathrm{align}}^{\mathcal I},
\qquad
\{\Delta_A^\ell\}_{A\subseteq V},
\qquad
 k_\varepsilon,
\qquad
\delta_A^Z.
\]
When the number of coalitions is too large for exhaustive evaluation, state the
sampling design, include every singleton and every leave-one-out coalition,
and report the most favorable coalition found for the reducibility null. The
claim must not be based only on average coalition performance.

\subtopic{74.6 Causal and reentry tests}

Report
\[
\eta_{\mathrm{causal}}^\nu,
\qquad
\Delta_{\mathrm{specific}}(r),
\qquad
\delta_{\mathrm{med}},
\qquad
\eta_{\mathrm{rec}}(\tau),
\qquad
\eta_{\mathrm{loop}}(\tau).
\]
Every causal estimate must be paired with:
\begin{enumerate}
\item a concrete mechanism-level implementation;
\item a manipulation check;
\item matched background variables;
\item declared direct paths that were blocked or preserved;
\item intervention magnitude and cost;
\item an off-workspace matched intervention; and
\item the downstream outcome window.
\end{enumerate}

\subtopic{74.7 Persistence, reconstruction, and replacement tests}

Report
\[
d_{\mathrm{id}}(W_t,W_{t+\tau}),
\qquad
\Rec_k,
\qquad
 e_t,
\qquad
\alpha^{(k)},
\qquad
\mu_t^{\mathrm{keep}}.
\]
The report must separate:
\begin{itemize}
\item subspace tracking from state tracking;
\item coherence from non-collapse;
\item reconstruction from mere prediction;
\item repair from successor-workspace creation; and
\item structural replacement from causal and reconstructive replacement.
\end{itemize}

\subtopic{74.8 Formation-provenance tests}

Report
\[
\mathbf E_{\mathrm{form}},
\qquad
\delta^{\mathrm{scaf}},
\qquad
\delta^{\mathrm{leak}},
\qquad
\Delta_q^{\mathrm{hist}}.
\]
Include the external scaffold, information available to the learner, history
shuffles, cross-scaffold retraining, and every route by which certification
labels or target boundaries could have entered formation.

\subtopic{74.9 Learned-structure and boundary tests}

Report
\[
\|\widehat H-H_*\|,
\qquad
\Delta_J,
\qquad
\Delta_t^B,
\qquad
\COut_\tau(\widehat S),
\qquad
\CIn_\tau(\widehat S).
\]
Boundary reports must include the admissible candidate family, full-system and
empty-set controls, singleton and centralized controls, score components,
margin to the runner-up, and temporal stability under held-out environment
drift.

\subtopic{74.10 Goal and counterfactual-policy tests}

Report the agency coordinates
\[
\eta_G(\tau),
\qquad
\delta_G^{\mathrm{cue}}(\tau),
\qquad
\eta_{G\to A},
\qquad
\delta_Q,
\qquad
\eta_{V\to A},
\qquad
\eta_{A\to Y},
\qquad
 k_{\mathrm{ctl},\varepsilon}.
\]
The required experiments include:
\begin{enumerate}
\item cue-controlled goal interventions;
\item goal-channel ablation and permutation;
\item counterfactual-value calibration against action interventions;
\item value-label permutation among feasible actions;
\item critic or value-channel blocking;
\item action replay under held-out perturbations;
\item policy regret or value-gap estimation; and
\item coalition reconstruction of the complete control state.
\end{enumerate}

\subtopic{74.11 Damage and repair tests}

For every damage family, report
\[
\eta_{D\to M},
\qquad
\eta_{M\to R},
\qquad
\eta_{R\to v},
\qquad
\Delta_{\mathrm{repair}}^{\mathrm{spec}},
\qquad
\delta_{\mathrm{rep-id}},
\qquad
\sup_t e_t.
\]
The repeated-disturbance experiment must contain:
\begin{enumerate}
\item a no-damage baseline;
\item damage with repair disabled;
\item damage with the intact repair loop;
\item monitor ablation;
\item cost-matched random repair;
\item off-workspace repair;
\item repeated and compound damage; and
\item a held-out disturbance distribution.
\end{enumerate}

\topic{75}{Required Null Families}

The core null panel is:
\[
\begin{aligned}
N_1&=\text{shuffled transfer maps with validated local matching},\\
N_2&=\text{shuffled cross-node temporal correspondence},\\
N_3&=\text{matched-capacity local memory},\\
N_4&=\text{matched-size irrelevant artifact},\\
N_5&=\text{participant ablation},\\
N_6&=\text{artifact ablation},\\
N_7&=\text{workspace implementation intervention},\\
N_8&=\text{matched off-workspace intervention},\\
N_9&=\text{representation replacement},\\
N_{10}&=\text{context reset and reconstruction},\\
N_{11}&=\text{matched formation-history shuffle},\\
N_{12}&=\text{cross-scaffold re-encoding and retraining},\\
N_{13}&=\text{boundary-label and certification leakage audit},\\
N_{14}&=\text{empty, full, singleton, and centralized boundary nulls},\\
N_{15}&=\text{held-out transfer and membership re-estimation},\\
N_{16}&=\text{boundary damage, use, and repair intervention},\\
N_{17}&=\text{goal-code permutation with cues and actions matched},\\
N_{18}&=\text{persistent external-cue rewrite with internal retention cut},\\
N_{19}&=\text{counterfactual-value label permutation},\\
N_{20}&=\text{critic or value-channel ablation},\\
N_{21}&=\text{open-loop action replay under held-out perturbation},\\
N_{22}&=\text{monitor-blind cost-matched random repair},\\
N_{23}&=\text{off-workspace repair with matched intervention norm},\\
N_{24}&=\text{identity-breaking reconstruction with matched task score},\\
N_{25}&=\text{merge control for components sharing an uninterruptible controller},\\
N_{26}&=\text{secret-sharing split of a known centralized controller},\\
N_{27}&=\text{hidden shared-memory and common-cue exposure test},\\
N_{28}&=\text{frozen certificate sweep across the admissible decomposition family},\\
N_{29}&=\text{post hoc target-selection audit},\\
N_{30}&=\text{horizon cherry-picking audit},\\
N_{31}&=\text{loss and utility substitution audit},\\
N_{32}&=\text{restricted intervention-family audit},\\
N_{33}&=\text{threshold and gate-set tuning audit},\\
N_{34}&=\text{frozen certificate sweep across the admissible specification family}.
\end{aligned}
\]

\subtopic{75.1 Null validity record}

Each observational null must state:
\begin{itemize}
\item the local distributions it preserves;
\item the relational dependency it destroys;
\item the distance between intact and null local marginals;
\item the capacity, parameter count, and information budget; and
\item whether null fitting used the same data budget as the intact model.
\end{itemize}
Each causal null must state:
\begin{itemize}
\item the manipulated mechanism;
\item the matched background and intervention cost;
\item the intended causal path destroyed or preserved;
\item the manipulation check; and
\item unintended direct effects detected during validation.
\end{itemize}
A failed manipulation check invalidates the null. It does not count as evidence
for the workspace claim.

\subtopic{75.2 Claim-specific null subsets}

A successful claim requires
\[
\Pass(N_1,\ldots,N_{28})
\]
only for the subset relevant to its stated level. The subset must be frozen
before certification. Failed applicable nulls, invalid nulls, and unmeasured
nulls are retained separately in the final report. They must not be merged
into a generic ``not significant'' category.

\topic{76}{Certification Vector and Complete Reporting Templates}

\subtopic{76.1 Extended certification vector}

Use a vector rather than one aggregate score:
\[
\boxed{
\begin{aligned}
\mathcal C_{\mathrm{full}}(W)
=\bigl(&
 d_\varepsilon,
 g_\varepsilon,
 \delta_{\mathrm{coh}},
 \delta_{\mathrm{align}}^{\mathcal I},
 \Delta_{\mathrm{joint}}^\ell,
 k_\varepsilon,\\
&\boldsymbol\Delta_{\mathrm{tr}}(\mathcal Q),
 \eta_{\mathrm{causal}}^\nu,
 \Delta_{\mathrm{specific}}(r),
 \eta_{\mathrm{loop}}(\tau),\\
&d_{\mathrm{id}},
 \Rec_k,
 \underline{\mathbf m}_{\mathfrak D},
 \boldsymbol\Delta_{\mathrm{dec}},
 \mathbf E_{\mathrm{form}},
 \delta^{\mathrm{scaf}},
 \delta^{\mathrm{leak}},\\
&\Delta_J,
 \Delta_t^B,
 \COut_\tau(\widehat S),
 \CIn_\tau(\widehat S),\\
&\eta_G(\tau),
 \delta_G^{\mathrm{cue}}(\tau),
 \eta_{G\to A},
 \delta_Q,
 \eta_{V\to A},
 \eta_{A\to Y},\\
&\eta_{D\to M},
 \eta_{M\to R},
 \eta_{R\to v},
 \Delta_{\mathrm{repair}}^{\mathrm{spec}},
 \delta_{\mathrm{rep-id}},
 \sup_t e_t
\bigr).
\end{aligned}
}
\]
A claim passes only when every required coordinate lies in its predeclared
admissible region. No weighted average may compensate for a failed required
gate.

\subtopic{76.2 Claim card template}

\begin{longtable}{@{}>{\raggedright\arraybackslash}p{0.29\textwidth}>{\raggedright\arraybackslash}p{0.65\textwidth}@{}}
\toprule
Field & Required entry \\
\midrule
\endfirsthead
\toprule
Field & Required entry \\
\midrule
\endhead
Claim identifier & Unique versioned identifier tied to code and result
artifacts. \\
Claim level & Archive, shared representation, operational workspace,
relational workspace, persistent workspace, self-maintaining workspace, or
workspace agent. \\
Reference decomposition & Nodes, artifacts, channels, local-access maps, and coalition universe. \\
Admissible decomposition family & Plausible merges, splits, shared-resource assignments, exclusions, and audit rationale. \\
Target and horizon & \(Y^+\), loss, utility, prediction horizon, control
horizon, and maintenance horizon. \\
Environment scope & Training, certification, and stress-test environment
families. \\
Workspace extractor & Frozen construction of \(D_H\), \(L_H\), spectral
threshold, and state readout. \\
Goal and policy & Goal codes, utility normalization, action set, and policy
class. \\
Boundary & Candidate family, score, admissibility constraints, and selected
boundary. \\
Interventions & Feasible workspace, goal, action, damage, and repair
implementations. \\
Null panel & Applicable \(N_j\), preservation targets, and fit budgets. \\
Thresholds & Every gate direction, numerical cutoff, and uncertainty rule. \\
Data split & Formation, calibration, certification, stress, and cross-fitting
logic. \\
Stopping rule & Number of seeds, runs, interventions, and any sequential
stopping criterion. \\
Status & Confirmatory, exploratory, replication, or audit. \\
\bottomrule
\end{longtable}

\subtopic{76.3 Gate-estimate table template}

For every required coordinate, include one row with the following fields:

\begin{longtable}{@{}>{\raggedright\arraybackslash}p{0.20\textwidth}>{\raggedright\arraybackslash}p{0.73\textwidth}@{}}
\toprule
Field & Required entry \\
\midrule
\endfirsthead
\toprule
Field & Required entry \\
\midrule
\endhead
Gate & Formal symbol and verbal interpretation. \\
Estimator & Exact statistic, normalization, and software implementation. \\
Evaluation unit & Episode, participant, graph, seed, environment, or other
independent unit. \\
Point estimate & Estimate on the frozen certification split. \\
Uncertainty & Interval, bound, or posterior region with declared coverage or
credibility. \\
Threshold & Predeclared pass region. \\
Manipulation status & Passed, failed, or not applicable. \\
Decision & Pass, fail, invalid, underpowered, or unmeasured. \\
Sensitivity & Result under reasonable alternative normalizations and
resampling choices. \\
Artifact reference & File, table, raw-result key, and code commit. \\
\bottomrule
\end{longtable}

\subtopic{76.4 Intervention ledger template}

\begin{longtable}{@{}p{0.20\textwidth}p{0.73\textwidth}@{}}
\toprule
Field & Required entry \\
\midrule
\endfirsthead
\toprule
Field & Required entry \\
\midrule
\endhead
Intervention ID & Unique identifier and randomization seed. \\
Target mechanism & Workspace coordinate, goal code, value channel, action,
boundary, monitor, or repair channel. \\
Implementation & Exact code path, hook, artifact edit, or mechanism change. \\
Dose & Norm, duration, timing, and support of the intervention. \\
Background match & Variables fixed, blocked, randomized, or balanced. \\
Excluded paths & Direct paths intentionally removed from the causal contrast. \\
Manipulation check & Evidence that the intended mechanism changed by the
specified amount. \\
Collateral effects & Changes outside the target mechanism. \\
Outcome & Downstream variable and evaluation window. \\
Effect estimate & Point estimate, uncertainty, and sign convention. \\
Cost & Compute, latency, information, action, or structural cost. \\
Validity status & Valid, invalid, ambiguous, or unmeasured. \\
\bottomrule
\end{longtable}

\subtopic{76.5 Null ledger template}

\begin{longtable}{@{}p{0.20\textwidth}p{0.73\textwidth}@{}}
\toprule
Field & Required entry \\
\midrule
\endfirsthead
\toprule
Field & Required entry \\
\midrule
\endhead
Null ID & One of \(N_1,\ldots,N_{34}\) or a declared extension. \\
Claim attacked & The exact explanation the null is designed to preserve or
destroy. \\
Matched quantities & Local marginals, capacity, training budget, action cost,
intervention norm, or other controlled quantities. \\
Destroyed relation & Edge alignment, temporal order, value ranking, goal
identity, monitor mediation, or another relational feature. \\
Fit procedure & Data, optimization, seeds, and hyperparameter budget. \\
Validity metric & Distance from the intended matched quantities. \\
Primary contrast & Intact minus null effect with uncertainty. \\
Disposition & Rejects null, fails to reject, invalid null, or underpowered. \\
Failure notes & Any mismatch or collateral mechanism change. \\
\bottomrule
\end{longtable}

\subtopic{76.6 Agency and maintenance report template}

The agency section must provide the following matrix:

\begin{center}
\small
\begin{tabularx}{\textwidth}{@{}p{0.22\textwidth}p{0.23\textwidth}X@{}}
\toprule
Component & Primary metric & Required contrast \\
\midrule
Goal semantics & Held-out \(Q_t^g\) ordering & Distinct goals reverse or alter
at least one declared action ordering. \\
Goal retention & \(\eta_G(\tau)\) & Goal intervention with future cues matched
or removed. \\
Goal use & \(\eta_{G\to A}\) & Goal intervention versus goal-label replay. \\
Counterfactual accuracy & \(\delta_Q\) & Internal value versus external action
interventions. \\
Value use & \(\eta_{V\to A}\) & Value-label permutation and value-channel
ablation. \\
Outcome control & \(\eta_{A\to Y}\) & Feasible action intervention. \\
Damage detection & \(\eta_{D\to M}\) & Damage versus matched no-damage. \\
Repair mediation & \(\eta_{M\to R}\) & Monitor intervention and monitor
ablation. \\
Restoration & \(\eta_{R\to v}\) & Repair versus matched no-repair. \\
Repair specificity & \(\Delta_{\mathrm{repair}}^{\mathrm{spec}}\) & Intended
repair versus cost-matched generic action. \\
Identity preservation & \(\delta_{\mathrm{rep-id}}\) & Pre-damage workspace
versus post-repair workspace. \\
Long-run maintenance & \(\sup_t e_t\) & Repeated held-out disturbances. \\
Relational location & \(k_{\mathrm{ctl},\varepsilon}\) & Full system versus
all single nodes and matched local nulls. \\
\bottomrule
\end{tabularx}
\end{center}

\subtopic{76.7 Uncertainty and multiplicity}

The independent resampling unit must match the data-generating process. Token,
time-step, or edge resampling is invalid when dependence is organized at the
episode, participant, graph, or environment level. Recommended procedures
include:
\begin{itemize}
\item environment- or episode-level bootstrap;
\item repeated independent training seeds;
\item cluster-robust intervals for repeated observations;
\item permutation tests that preserve the declared local structure;
\item concentration bounds for uniform finite candidate families; and
\item stress-test intervals over held-out environment families.
\end{itemize}

Pass decisions should use one-sided uncertainty bounds. For a lower-bounded
gate \(\theta\geq\tau\), pass only when the lower confidence bound exceeds
\(\tau\). For an upper-bounded defect \(\delta\leq\varepsilon\), pass only when
the upper confidence bound is below \(\varepsilon\).

When a conclusion ranges over \(\mathfrak S\), the uncertainty rule must be
uniform over every gate and specification included in the claim. Report the
pointwise interval, the family-wide lower envelope, the covering or
multiplicity correction, and the identification status. A narrow pointwise
interval does not establish a robust result when model or specification
uncertainty remains wide.

\subtopic{76.8 \(T_{24}\): Gatewise certification validity }

\begin{theorem}[Intersection-gate false-certification control]
\label{thm:gatewise-validity}
Let a certificate require all gates \(G_1,\ldots,G_m\) to pass. For each gate
\(j\), let \(H_{0j}\) denote the null that the required property is absent, and
assume the predeclared test satisfies
\[
\sup_{P\in H_{0j}}
\mathbb P_P(G_j\text{ passes})
\leq
\alpha_j.
\]
The global null is
\[
H_0
=
\bigcup_{j=1}^m H_{0j},
\]
meaning that at least one required property is absent. If the certificate is
declared only when every gate passes, then
\[
\boxed{
\sup_{P\in H_0}
\mathbb P_P(\text{certificate passes})
\leq
\max_{1\leq j\leq m}\alpha_j.
}
\]
No independence assumption among gate tests is required.
\end{theorem}

\begin{proof}
For any \(P\in H_0\), choose an index \(j_P\) such that \(P\in H_{0j_P}\).
The event that every gate passes is contained in the event that gate \(j_P\)
passes. Hence
\[
\mathbb P_P(\text{certificate passes})
\leq
\mathbb P_P(G_{j_P}\text{ passes})
\leq
\alpha_{j_P}
\leq
\max_j\alpha_j.
\]
Taking the supremum over \(P\in H_0\) proves the claim.
\end{proof}

The theorem is an intersection-union guarantee. It applies only when the gate
family, thresholds, and tests are fixed before certification and each gate test
is uniformly valid for its null. Adaptive gate deletion, selective reporting,
or repeated reuse of the certification split invalidates the bound.

\subtopic{76.9 Failure taxonomy}

Every gate and null receives exactly one primary status:
\begin{description}
\item[Pass.] The frozen estimate and uncertainty rule satisfy the threshold.
\item[Fail.] The gate was validly measured and did not satisfy the threshold.
\item[Invalid intervention.] The manipulation check or background match
 failed.
\item[Invalid null.] The null did not preserve its required local quantities
 or exceeded its fit budget.
\item[Underpowered.] The uncertainty interval crosses the threshold by more
 than the declared indifference region.
\item[Unmeasured.] No valid estimate was obtained.
\item[Leakage.] Formation or model selection used certification information.
\item[Exploratory.] The metric, threshold, or analysis was selected after
 inspecting certification outcomes.
\end{description}
Only \emph{Pass} counts toward the formal certificate. ``Underpowered'' and
``unmeasured'' are not failures of the underlying theory, but they are failures
to establish the corresponding empirical claim.

\subtopic{76.10 Machine-readable result record}

The machine-readable record should contain at least the following top-level
objects:
\[
\boxed{
\begin{aligned}
\mathsf{ResultRecord}
=\{&
\mathsf{claim},
\mathsf{specification},
\mathsf{specification\_family},
\mathsf{decomposition},
\mathsf{data\_splits},
\mathsf{formation},\\
&\mathsf{workspace\_extractor},
\mathsf{goals},
\mathsf{actions},
\mathsf{interventions},\\
&\mathsf{nulls},
\mathsf{gates},
\mathsf{identified\_sets},
\mathsf{specification\_envelopes},
\mathsf{uncertainty},
\mathsf{failures},
\mathsf{artifacts}
\}.
\end{aligned}
}
\]
Each gate object records its symbol, estimator, direction, threshold, point
estimate, interval, split, seed set, status, and artifact reference. Each
intervention and null object records its manipulation check and validity
status. The schema must preserve failed and invalid entries rather than only
the passing subset.

\subtopic{76.11 Minimum reproducibility package}

A complete certification package contains:
\begin{enumerate}
\item the frozen claim card;
\item source code and environment lock file;
\item exact data or deterministic data-generation procedure;
\item formation, calibration, certification, and stress-test manifests;
\item learned transfer, workspace, policy, monitor, and repair parameters;
\item raw per-run and per-intervention results;
\item the intervention and null ledgers;
\item the complete gate table, including failures;
\item figures and tables regenerated from raw results;
\item random seeds and hardware or numerical-precision information;
\item checksums for every included artifact; and
\item a single verification command that reproduces package integrity and the
 reported summary tables.
\end{enumerate}

\subtopic{76.12 Decision rule}

Let \(\mathcal J(L)\) be the required gate set for claim level \(L\). Then
\[
\boxed{
\operatorname{Certify}_L(W;\varsigma)
=
\bigwedge_{j\in\mathcal J(L,\varsigma)}
\operatorname{Pass}_j(W;\varsigma).
}
\]
For a workspace-agent claim, the required set includes every structural,
operational, relational, persistent, goal, counterfactual-control, and
self-maintenance gate declared in Parts~IX and XVI. A specification-robust
claim additionally requires the uniform rule of Part~XVIII across
\(\mathfrak S\). Failure of an earlier level
prevents certification of a later level, even when later-level behavior appears
strong on a task benchmark.
